\section{Related Work}
\label{sec:related_work}

The method above connects two themes that are often treated separately: how to choose a stable teacher target, and how to assign credit along an on-policy trajectory. We organize related work around this split. Prior work provides ingredients for smoother supervision, reference-regularized policy optimization, and return-based credit assignment; \ours{} combines these ingredients into a single distillation procedure.

\paragraph{On-policy distillation and self-distillation.} On-policy distillation trains the student on its own generated trajectories, using dense
teacher feedback at each token. This avoids the exposure-bias mismatch of purely offline
distillation, but introduces a new challenge: the teacher must provide reliable supervision
on student-generated prefixes, which may be outside the teacher's typical distribution \citep{li2026rethinking, jang2026stable}.
Recent self-distillation methods further instantiate the teacher as the same model with
access to privileged information, such as ground-truth solutions or feedback \citep{zhao2026self, liu2026self, yang2026self, shen2026anti, zhang2026opsdl, hubotter2026reinforcement, he2026self}. \ours{} belongs to this line of work, but differs by replacing the fixed teacher target with a scheduled
student--teacher interpolant and by correcting the myopic local token-KL gradient with a
Look-Ahead return.

\paragraph{RLHF and reference-regularized policy optimization.}
Our theoretical formulation is closely related to KL-regularized reinforcement learning and
RLHF, where a policy is optimized to maximize reward while staying close to a reference
model \citep{ziegler2019fine, ouyang2022training, schulman2017proximal}. DPO shows that the optimal policy of
a KL-regularized preference objective can be written in closed form, revealing a direct
connection between rewards, reference policies, and optimal distributions. We use a similar
perspective to reinterpret on-policy distillation: the teacher-to-reference log-probability
ratio acts as a reward, and standard OPD is recovered as the $\beta=1$ special case. This
view leads directly to the geometric interpolant between the reference student and the
privileged teacher.

\paragraph{Interpolation, curricula, and teacher-guided learning.}
Curriculum learning aims to make optimization smoother by gradually changing the
difficulty or source of supervision \citep{bengio2009curriculum}. Related ideas also
appear in sequence prediction and imitation learning: scheduled sampling gradually moves
training from teacher-forced inputs toward model-generated inputs
\citep{bengio2015scheduled}, while DAgger addresses distribution shift by querying expert
supervision on learner-induced states \citep{ross2011reduction}. More recent work has
explored student--teacher mixture reformulations to bridge the distributional gap in
distillation \citep{jang2026stable}, as well as teacher-guided or
student--teacher cooperative sampling for improving training rollouts
\citep{xu2025speculative,shenfeld2023tgrl}. Our logit-interpolant target can be viewed as
a distributional curriculum: early targets remain close to the student, while later targets
move toward the privileged teacher. Unlike heuristic interpolation or guidance used only
for sampling, our interpolant is derived as the optimal policy of a reference-regularized
distillation objective. This provides a principled target path from student to teacher and,
when combined with return-to-go credit assignment, yields a stable and future-aware
on-policy distillation algorithm.

\paragraph{Policy gradients and credit assignment.}
The return-to-go estimator used in $\beta$-OPSD is closely related to classical policy-gradient methods such as
REINFORCE \citep{williams1992simple}, where unbiased gradient estimation assigns each action its
return-to-go rather than only its instantaneous reward. In OPD, the analogous per-step
signal is the teacher--student log-ratio along a generated trajectory. Prior work has explored
using local token-level signals for stable distillation \citep{li2026rethinking, armandpour2026unmasking, shen2026anti, jang2026stable, kim2026does}, and related derivations have
also identified unbiased sequence-level gradient estimators for reverse-KL objectives
\citep{agrawal2026reinforcement, liu2026self}. Future-aware credit assignment has also been studied in GRPO-style optimization
\citep{schulman2015high, sutton1998reinforcement}. Our work builds on this perspective and shows that, for on-policy distillation, the
standard local token-KL update is myopic because it ignores the future distributional
mismatch induced by earlier tokens. \ours{} applies return-to-go credit assignment directly
to teacher-guided distillation: the exact undiscounted Look-Ahead estimator recovers an
unbiased sequence-level KL gradient under student on-policy sampling, while discounted
implementations provide a practical stability--variance trade-off for long generations.